\documentclass[11pt]{article}
\usepackage[margin=1in]{geometry}
\usepackage{amsmath,amssymb,amsthm,mathtools}
\usepackage{graphicx}
\usepackage{microtype}
\usepackage[colorlinks=true,linkcolor=blue,citecolor=blue,urlcolor=blue]{hyperref}

\newtheorem{theorem}{Theorem}
\newtheorem{corollary}{Corollary}

\newcommand{\R}{\mathbb R}
\newcommand{\C}{\mathbb C}
\newcommand{\B}{\mathcal B}
\newcommand{\ip}[2]{\langle #1,#2\rangle}

\title{The Parameter-Precompactness Criterion for\\
Strong Bohr Approximation}
\author{Run \texttt{fa\_banach\_001}, agent \texttt{agent\_lane\_05}}
\date{2026-08-09}

\begin{document}
\maketitle

\begin{abstract}
We give a complete parameter-side answer on the canonical domain $I=\R^n$ to
the converse question following Proposition 2.24 of
Ch\'avez--Khalil--Kosti\'c--Pinto.  For a family $\B$ of subsets of a Banach
space $X$, every continuous Bohr $\B$-almost periodic function on
$\R^n\times X$ is strongly $\B$-almost periodic if and only if every member
of $\B$ is relatively compact.  Thus for the family of all bounded subsets
the universal implication holds exactly when $X$ is finite-dimensional.  The
same characterization holds on the full-dimensional convex polyhedral cones
used in the source.  In contrast, when $n\geq2$, diagonal almost periods alone
do not imply strong almost periodicity even for a bounded scalar function with
no parameter variable.
\end{abstract}

\section*{Source Question}

After Proposition 2.24(iv)--(v) of \cite{CKKP}, the authors ask:
\begin{quote}
Suppose that $\varnothing\neq I\subseteq\R^n$ and
$F:I\times X\to Y$ is Bohr $\mathcal B$-almost periodic (or Bohr
$(\mathcal B,I\cap\Delta_n)$-almost periodic).  What conditions ensure the
strong $\mathcal B$-almost periodicity of $F$?
\end{quote}

\begin{center}
\includegraphics[width=0.94\textwidth]{figures/open_problem_crop.png}
\end{center}

For $I=\R^n$, Bohr $\B$-almost periodicity means that for every
$B\in\B$ and $\varepsilon>0$ there is $L>0$ such that every ball of radius
$L$ contains a translation $\tau$ satisfying
\[
 \sup_{t\in\R^n,\,x\in B}
 \|F(t+\tau,x)-F(t,x)\|_Y\leq\varepsilon.
\]
Strong $\B$-almost periodicity means that, for every $B\in\B$, $F$ is the
uniform limit on $\R^n\times B$ of polynomials
\[
 P(t,x)=\sum_{j=1}^r e^{i\ip{\lambda_j}{t}}c_j(x),
 \qquad \lambda_j\in\R^n,\quad c_j\in C(X;Y).
\]

\section*{Proof Intuition}

There are two independent obstructions.  On the parameter side, compactness
turns $F$ into the ordinary Banach-valued Bohr function
$t\mapsto F(t,\cdot)|_K\in C(K;Y)$.  Its vector-valued trigonometric
approximants have coefficients in $C(K;Y)$, and Dugundji's theorem extends
those coefficients from $K$ to all of $X$.

This is sharp.  Any non-precompact parameter set contains a uniformly
separated sequence.  Disjoint bumps around that sequence can carry the
successive modes $e^{imt_1}$.  The resulting function has the exact common
period $2\pi$ but no finite-frequency polynomial can approximate all the
modes: a character mean gives error at least one.

Diagonal translations introduce a different, purely temporal obstruction.
They do not see the transverse coordinate $t_1-t_2$.  The bounded function
$e^{i(t_1-t_2)^2}$ is therefore exactly invariant under every diagonal shift,
but it cannot be a uniform limit of trigonometric polynomials because its
restriction $e^{is^2}$ is not uniformly continuous.

\section*{Exact Euclidean-Domain Characterization}

\begin{theorem}[parameter-precompactness criterion]\label{thm:main}
Let $n\geq1$, let $X$ be a complex Banach space, and let $\B$ be a nonempty
family of nonempty subsets of $X$.  The following are equivalent.
\begin{enumerate}
\item Every $B\in\B$ is relatively compact in $X$.
\item For every complex Banach space $Y$, every continuous Bohr
$\B$-almost periodic function $F:\R^n\times X\to Y$ is strongly
$\B$-almost periodic.
\item Statement \textup{(2)} holds for scalar-valued functions $Y=\C$.
\end{enumerate}
When these conditions hold and the members of $\B$ are bounded, Bohr and
strong $\B$-almost periodicity are equivalent.
\end{theorem}

\begin{proof}
Assume (1), fix $B\in\B$, and put $K=\overline B$.  Then $K$ is compact.
For
\[
 G_K(t)=F(t,\cdot)|_K\in C(K;Y),\qquad t\in\R^n,
\]
the map $G_K:\R^n\to C(K;Y)$ is norm-continuous.  Indeed, joint continuity
of $F$ at the points of $\{t_0\}\times K$, followed by a finite-cover
argument on $K$, gives
\[
 \sup_{x\in K}\|F(t,x)-F(t_0,x)\|_Y\longrightarrow0
 \quad(t\to t_0).
\]
For each fixed $t,\tau$, continuity in $x$ and density of $B$ in $K$ give
\[
 \sup_{x\in B}\|F(t+\tau,x)-F(t,x)\|_Y
 =\sup_{x\in K}\|F(t+\tau,x)-F(t,x)\|_Y.
\]
Hence the Bohr $\B$ estimate for $B$ says exactly that $G_K$ is a
$C(K;Y)$-valued Bohr almost periodic function.

The Banach-valued Bohr approximation theorem recalled in Section 1.1 of
\cite{CKKP} gives trigonometric polynomials
\[
 Q_q(t)=\sum_{j=1}^{r_q}e^{i\ip{\lambda_{q,j}}{t}}a_{q,j},
 \qquad a_{q,j}\in C(K;Y),
\]
such that
\[
 \sup_{t\in\R^n}\|G_K(t)-Q_q(t)\|_{C(K;Y)}\longrightarrow0.
\]
The set $K$ is closed in the metric space $X$, and $Y$ is locally convex.
Dugundji's extension theorem \cite{Dugundji} extends every $a_{q,j}$ to a
continuous map $c_{q,j}:X\to Y$.  Therefore
\[
 P_q^B(t,x)=\sum_{j=1}^{r_q}
 e^{i\ip{\lambda_{q,j}}{t}}c_{q,j}(x)
\]
has the exact form required in the source and converges uniformly to $F$ on
$\R^n\times K$, hence on $\R^n\times B$.  Since $B$ was arbitrary, $F$ is
strongly $\B$-almost periodic.  This proves (1)$\Rightarrow$(2), while
(2)$\Rightarrow$(3) is immediate.

We prove the contrapositive of (3)$\Rightarrow$(1).  Suppose that some
$B_0\in\B$ is not relatively compact.  Since $X$ is complete, $B_0$ is not
totally bounded.  Thus there are $\delta>0$ and points $x_m\in B_0$ with
\[
 \|x_m-x_k\|\geq\delta\qquad(m\neq k).
\]
Define
\[
 \phi_m(x)=\max\{0,1-3\|x-x_m\|/\delta\}.
\]
The supports lie in the radius-$\delta/3$ balls around the $x_m$ and are
pairwise separated.  The family $(\phi_m)$ is locally finite, and at most one
term is nonzero at any point.  Set
\begin{equation}\label{eq:F}
 F(t,x)=\sum_{m=1}^{\infty}\phi_m(x)e^{imt_1}.
\end{equation}
Then $F$ is jointly continuous and $|F|\leq1$.  Every
$\tau\in\R^n$ with $\tau_1\in2\pi\mathbb Z$ is an exact period.  These
periods are relatively dense: one changes the first coordinate of any point
to its nearest multiple of $2\pi$, moving by at most $\pi$.  Thus $F$ is
Bohr $\B$-almost periodic for every family $\B$.

Let
\[
 P(t,x)=\sum_{j=1}^{r}e^{i\ip{\lambda_j}{t}}c_j(x)
\]
be any allowed trigonometric polynomial.  Choose $m$ such that
$me_1\notin\{\lambda_1,\ldots,\lambda_r\}$.  Since
$F(t,x_m)=e^{imt_1}$, multiplication by $e^{-imt_1}$ and averaging over
$Q_T=[-T,T]^n$ gives
\begin{align*}
 &\frac1{(2T)^n}\int_{Q_T}e^{-imt_1}
 \bigl(F(t,x_m)-P(t,x_m)\bigr)\,dt\\
 &\qquad=1-\sum_{j=1}^{r}c_j(x_m)
 \frac1{(2T)^n}\int_{Q_T}
 e^{i\ip{\lambda_j-me_1}{t}}\,dt\longrightarrow1.
\end{align*}
Every nontrivial character has large-cube mean zero.  The absolute value of
the left side is at most the uniform error on $\R^n\times B_0$; hence
\[
 \sup_{t\in\R^n,\,x\in B_0}|F(t,x)-P(t,x)|\geq1.
\]
No strong approximation exists, contradicting (3).  This proves the
equivalence.

Finally, if the members of $\B$ are bounded, strong $\B$-almost periodicity
implies Bohr $\B$-almost periodicity by Proposition 2.24(iv) of \cite{CKKP}
(the domain is $\R^n$).  The last assertion follows.
\end{proof}

\begin{corollary}\label{cor:standard}
For the family of all compact subsets of $X$, Bohr and strong
$\B$-almost periodicity coincide on $\R^n$.  For the family of all bounded
subsets, the universal implication from Bohr to strong $\B$-almost
periodicity holds if and only if $X$ is finite-dimensional.
\end{corollary}

\begin{proof}
Compact sets are relatively compact.  All bounded subsets of a Banach space
are relatively compact exactly in finite dimension.
\end{proof}

\section*{Convex Polyhedral Domains}

The same criterion applies to the source's principal non-group domains.

\begin{theorem}\label{thm:cone}
Let $(v_1,\ldots,v_n)$ be a basis of $\R^n$ and
\[
 I=\left\{\sum_{j=1}^n a_jv_j:a_j\geq0\right\}.
\]
For a family $\B$ of subsets of a Banach space $X$, every continuous Bohr
$\B$-almost periodic map $F:I\times X\to Y$ is strongly $\B$-almost
periodic for every Banach target $Y$ if and only if every $B\in\B$ is
relatively compact.
\end{theorem}

\begin{proof}
For sufficiency, fix $B$ and $K=\overline B$.  As above,
$G_K(t)=F(t,\cdot)|_K$ is a continuous Bohr almost periodic map
$I\to C(K;Y)$.  The convex-polyhedral consequence of Theorem 2.36 in
\cite{CKKP}, applied to the Banach target $C(K;Y)$, uniformly approximates
$G_K$ on $I$ by $C(K;Y)$-valued trigonometric polynomials.  Dugundji
extension of their coefficients gives the required strong approximants.

For necessity, let $(x_m)$ and $(\phi_m)$ be the separated sequence and
bumps from the proof of Theorem \ref{thm:main}.  Let $\ell_1$ be the dual
coordinate with $\ell_1(v_1)=1$ and $\ell_1(v_j)=0$ for $j>1$, and put
\[
 F(t,x)=\sum_{m\geq1}\phi_m(x)e^{im\ell_1(t)},\qquad t\in I.
\]
This function is bounded and continuous.  Its exact periods in $I$ include
all $\tau$ with $\ell_1(\tau)\in2\pi\mathbb Z$.  They are relatively dense
in $I$: if $t^0=\sum a_jv_j$, replace $a_1$ by its nearest nonnegative
multiple of $2\pi$, changing $t^0$ by at most $\pi\|v_1\|$.

If a finite trigonometric polynomial approximated $F$ on $I\times B_0$,
restricting at $x_m$ to the ray $t=sv_1$, $s\geq0$, would approximate
$e^{ims}$ by finitely many characters $e^{i\ip{\lambda_j}{v_1}s}$.  Choose
$m$ outside that finite frequency set and average over $[0,T]$ after
multiplication by $e^{-ims}$.  The same character-mean argument gives error
at least one.  Thus non-precompact $B_0$ again contradicts the universal
implication.
\end{proof}

\section*{Diagonal Periods Do Not Suffice}

\begin{theorem}\label{thm:diagonal}
Let $n\geq2$.  Even with no effective parameter variable, there is a bounded
continuous scalar function on $\R^n$ that is Bohr
$(\B,\Delta_n)$-almost periodic but is not strongly $\B$-almost periodic.
Consequently no condition on the parameter family alone can reverse
Proposition 2.24(v).
\end{theorem}

\begin{proof}
Take $X=\{0\}$, $\B=\{\{0\}\}$, and
\[
 F(t)=e^{i(t_1-t_2)^2}.
\]
Every diagonal vector $(\tau,\ldots,\tau)$ is an exact period, so the
diagonal Bohr condition holds with zero error.

If $F$ were strongly almost periodic, then its restriction to the line
$t=(s,0,\ldots,0)$ would be a uniform limit on $\R$ of ordinary
trigonometric polynomials.  Such a limit is uniformly continuous.  But
$e^{is^2}$ is not uniformly continuous: for
\[
 a_k=\sqrt{2\pi k},\qquad b_k=\sqrt{2\pi k+\pi},
\]
we have $|a_k-b_k|\to0$ while $e^{ia_k^2}=1$ and $e^{ib_k^2}=-1$.
This contradiction proves the claim.
\end{proof}

\section*{Verification Notes and Scope}

Theorem \ref{thm:main} upgrades the earlier compact-versus-bounded example
to an exact criterion for every parameter family on $\R^n$.  Its necessity
uses only completeness of $X$ and the elementary equivalence between failure
of total boundedness and the existence of a separated sequence.  It is not
special to $\ell_2$.

Theorem \ref{thm:cone} covers the full-dimensional convex polyhedral cones
explicitly used after Theorem 2.36 in the source.  Arbitrary additive domains
$I\subsetneq\R^n$ may have additional temporal extension obstructions and are
not classified here.  Theorem \ref{thm:diagonal} completely rules out a
parameter-only converse for diagonal periods in dimensions at least two.

\section*{Novelty Check}

The run indexes and local arXiv source corpus were searched for
arXiv:2012.00543, the exact converse question, and combinations of
``strong $\mathcal B$-almost periodic'', ``Bohr $\mathcal B$-almost
periodic'', ``relatively compact'', and ``precompact''.  Bounded web searches
on 2026-08-09 checked the original paper and later work on metrical and
Stepanov variants.  They found the source and related generalizations, but no
later answer, parameter-precompactness iff criterion, or separated-frequency
counterexample matching the theorems above.  This is a bounded novelty check,
not a certification of originality.

\section*{Human-Review Recommendation}

Treat Theorems \ref{thm:main}--\ref{thm:diagonal} as a likely-valid complete
classification for $I=\R^n$, with a convex-polyhedral extension.  The key
review points are the $C(K;Y)$ norm-continuity step, Dugundji extension of
vector-valued coefficients, local finiteness of the separated bump family,
the character-mean lower bound, and the exact scope of the source's Theorem
2.36 used in Theorem \ref{thm:cone}.  No conditional lemma remains.

\begin{thebibliography}{9}

\bibitem{CKKP}
A. Ch\'avez, K. Khalil, M. Kosti\'c, and M. Pinto,
\emph{Multi-dimensional almost periodic type functions and applications},
arXiv:2012.00543 (2020).

\bibitem{Dugundji}
J. Dugundji,
\emph{An extension of Tietze's theorem},
Pacific Journal of Mathematics 1 (1951), no.~3, 353--367.

\end{thebibliography}

\end{document}
